
%Inicio del Preámbulo: Modifícalo bajo tu propio riesgo! Usa los paquetes que requieras para 
%implementar funcionalidades dentro del documento. Se comenta solo lo esencial, puedes buscar los
%paquetes y documentarte de su funcionalidad y puedes eliminarlos también si así lo deseas. 

\documentclass[letterpaper,12pt]{article} %Modifica el tipo de documento y el tamaño de la letra.
\usepackage[utf8]{inputenc} %Formato UTF-8 para caracteres especiales.
\usepackage[shortlabels]{enumitem}
\usepackage[spanish,mexico]{babel} 
\usepackage{amsmath,amssymb,amsfonts,latexsym,cancel}
\usepackage{hyperref}
\usepackage{wrapfig}
\usepackage[rflt]{floatflt}
\usepackage[pdftex]{graphicx}
\usepackage{fancyhdr} %Paquete para el header y el formato de la portada. No sugiero borrarlo!
\usepackage{float}
\usepackage{longtable,multirow,booktabs}
\usepackage{cite}
\usepackage{wrapfig}
\usepackage[square,numbers]{natbib}
\usepackage{multicol}
\usepackage{caption}
\usepackage[]{sidecap}
\usepackage{adjustbox}
\usepackage{parskip}
\usepackage{enumitem}
\usepackage{tikz}
\usepackage{lipsum}
\usepackage[]{xcolor}



%Fin de Préambulo


%Inicio formato de Página. Puedes establecer medidas de la página y modificar el header.

\textheight = 21cm %Medidas de la  página
\textwidth = 18cm  %Medidas de la página
\topmargin = -2cm  %Medidas de la página    
\oddsidemargin = -0.8cm %Medidas de la página
\pagestyle{fancy} %Diseño de la página


\fancyhf{}
\lhead{Facultad de Ingeniería}%%LeftHead
\chead{\includegraphics[width=1cm, height=1cm]{Imágenes/UNAM.png}}%%CenterHead
%\lfoot{USM}
\rhead{Asignatura Teórica o Laboratorio}%%RightHead

\setlength{\columnsep}{4mm}%Comandos para el formato de la página.
%\setlength{\parindent}{4em}%Sangría al comenzar un nuevo párrafo.
\setlength{\parindent}{0.5in}
%\setlength{\parindent}{4em}%Sangría al comenzar un nuevo párrafo.
\setlength{\parskip}{1em}%Distancia entre párrafos.
\renewcommand{\baselinestretch}{1.5}% Espacio entre línea y línea o interlineado.
\setlength{\headheight}{33pt}
\fancyfoot[CE,CO]{\LaTeX{}\\\thepage} %Logo de LaTeX y pie de página.

%Fin formato de Página

%Aquí inicia el documento. Date como magnate
\begin{document}

    %LaTeX te hace el índice automáticamente conforme añades secciones en tu documento.
    \input{portada.tex}
    \tableofcontents
    \newpage
%Inicio parte opcional. Esta parte la puedes quitar si deseas, es por si te piden formatos para
%evidencias de certificación de los laboratorios con números de cuenta o te piden abstracts en tus %trabajos.

\title{Nombre del curso o laboratorio \\\textbf{Título del reporte o trabajo.} \\ } 

\author{
  \normalsize{\texttt{123456789}} \and
  \normalsize{\texttt{123456789}}
  \and
 \normalsize{\texttt{123456789}}
  \and
  \normalsize{\texttt{123456789}}
  \and
  \normalsize{\texttt{123456789}}
}
\date{1 de enero, 2022}
\maketitle
\thispagestyle{fancy}

\begin{abstract}
En esta parte del reporte se coloca lo que se conoce como un \textit{abstract}, es decir, un resumen de la temática del reporte de laboratorio o artículo presentado que incluye los elementos o palabras clave que se retomarán durante el desarrollo del mismo.
\end{abstract}

%Fin parte opcional

\section{Introducción}

\noindent En esta sección se coloca la introducción a la temática abordada por primera vez en el abstract. Introduce los conceptos a tratar en un enfoque teórico y señala los elementos que serán empleados para el desarrollo experimental del reporte o trabajo. Nota que este párrafo no está indentado. Si quieres quitar el indentado manualmente, con noindent puedes hacerlo.

\section{Metodología de la práctica}

En esta sección se pueden introducir la forma de trabajo que se empleará para realizar las actividades experimentales a realizar en la práctica. Si no requieres esta sección, puedes eliminarla, o bien, modificarla para que se amolde a tus necesidades.

Si necesitas ecuaciones en una línea separada sin numerar, puedes usar esto:

$$s_f=s_0+v_0t+\frac{1}{2}at^2$$

o las puedes escribir en línea de esta forma: $v_f^2=v_i^2+2a\Delta s$ o $\xi$.

\section{Experimentos}

Aquí puedes colocar cada uno de los experimentos que realizaste para la elaboración de tu práctica, trabajo, o proyecto. Aquí tienes un ejemplo de como colocar una imagen. No te tienes que preocupar por enumerarlas, ¡LaTeX ya lo hace por ti! Las puedes citar en cualquier parte como [\ref{exemploLabel}].

\begin{figure}[h!]
    \centering
    %Modificando los parámetros de height puedes cambiar el tamaño de tu imagen.
    \includegraphics[height=9cm]{Imágenes/DummyExp.jpg}
    \caption{\textit{Simón profe, si le entendimos a la práctica.}}
    \label{exemploLabel}
    \end{figure}

\section{Resultados}

Aquí puedes poner todos los resultados que obtuviste, incluyendo tablas, gráficos, dibujos o esquemas. Aquí abajo tienes como ejemplo una tabla con LaTeX. Puedes crearlas con generadores online, o si se te complica, lo más fácil es hacerlas en Excel y solo capturar la imagen para incluirla como el ejemplo anterior. 

\begin{table}[ht]
\begin{center}
\caption{Todas las tablas necesitan un pie de foto en la parte superior. Fijate en la consistencia de los valores de las tablas. No olvides tus unidades.}
\label{table1} 
\begin{tabular}{ccc} %change to cc for 2 columns
\hline
\multicolumn{1}{c}{Distance, $d$ (km) } & \multicolumn{1}{c}{Voltage, $V\ (\pm 0.05$ V)} & \multicolumn{1}{c}{Current, $I$\ (mA $\pm 5$\%)}\\
\hline
1.2 $\pm$ 0.2 &  0.30 & 20 \\
1.6 $\pm$ 0.4 &  0.21 & 30 \\
2.5 $\pm$ 0.1 &  0.18 & 40 \\
5.9 $\pm$ 0.2 &  0.13 & 50 \\
\hline
\end{tabular}
\end{center}
\end{table}


\section{Conclusion}
Finalmente, aquí puedes poner tus conclusiones, sean individuales o por equipo. Si son por equipo, una forma fácil de hacerlo es con una lista itemize:

\begin{itemize}
    \item \textbf{Nombre 1:} si me gustó la práctica jsjs.
    
    \item \textbf{Nombre 1:} no me gustó la práctica :(
    
\end{itemize}

Por último, tienes las referencias que usaste. Si eres muy pro y no usaste nada, puedes comentar las siguientes tres líneas y no aparecerá esa sección. Agrega tus referencias en el archivo Libros.bib y citalas aquí.

            \nocite{*}
            \bibliographystyle{apalike}
            \bibliography{Libros.bib}


%Pequeño disclaimer por si usas fotos con derechos de autor que sacaste de Internet. Nunca está de más. 
\centering \small Los créditos de las fotografías pertenecen a sus respectivos autores. ©

\centering\vspace*{\fill} 
\end{document}

 %Fin del documento.